

% \begin{figure*}[t]
%     \centering
%     \hspace*{-0.2cm} % 核心负向位移指令

%     \includegraphics[width=1.0\linewidth]{top.png}
    
%     \caption{{\color{red}This figure illustrates the full-resolution training renderings of the ``bicycle'' scene in the mip-nerf 360 dataset, as well as the average GPU memory consumption and model sizes of 3DGS, Mip-Splatting, and the proposed method. The lower part includes two charts: the left-hand chart compares the model sizes, and the right-hand chart illustrates the GPU memory consumption, showing that Mip-Splatting has the highest usage and Ours has the lowest.}}
%     \label{fig:top}
% \end{figure*}



\begin{abstract}
% 3D Gaussian Splatting (3DGS) has significantly driven progress in real-time 3D scene reconstruction. However, its application in high-resolution reconstruction scenarios encounters critical memory scalability limitations. To address this, we propose Hierarchically Refined Gaussian Splatting (HRGS), a memory-efficient framework with hierarchical block optimization from coarse to fine. Specifically, we first derive a coarse global Gaussian representation from low-resolution data, then partition the scene into blocks and refine each block using high-resolution data. The scene partitioning includes two procedures: Gaussian partitioning and training data partitioning. For Gaussian partitioning, we contract irregular scenes into normalized bounded cubic shapes and employ uniform grid division to distribute computational tasks evenly among various blocks. For training data partitioning, only the observations that reside within their corresponding blocks or significantly contribute to the rendering results are retained. Alignment of trained Gaussians from adjacent blocks is achieved by guiding each block's refinement with the coarse global Gaussian prior, ensuring seamless fusion. To further enhance surface reconstruction quality, we incorporate normal priors from a pretrained model. Finally, our method enables high-quality and high-resolution 3D scene reconstruction even under constrained memory capacities. Extensive experimental results on three public benchmarks show that our approach achieves state-of-the-art performance in 3D reconstruction, covering both high-resolution novel view systhesis (NVS) and surface reconstruction.

3D Gaussian Splatting (3DGS) has achieved significant progress in real-time 3D scene reconstruction. However, its application in high-resolution reconstruction scenarios faces severe memory scalability bottlenecks. To address this issue, we propose Hierarchical  Gaussian Splatting (HRGS), a memory-efficient framework with hierarchical block-level optimization from coarse to fine. Specifically, we first derive a global, coarse Gaussian representation from low-resolution data; we then partition the scene into multiple blocks and refine each block using high-resolution data. Scene partitioning comprises two steps: Gaussian partitioning and training data partitioning. In Gaussian partitioning, we contract irregular scenes into a normalized, bounded cubic space and employ a uniform grid to evenly distribute computational tasks among blocks; in training data partitioning, we retain only those observations that lie within their corresponding blocks or make significant contributions to the rendering results. By guiding each block’s refinement with the global coarse Gaussian prior, we ensure alignment and seamless fusion of Gaussians across adjacent blocks. To reduce computational resource demands, we introduce an Importance-Driven Gaussian Pruning (IDGP) strategy: during each block’s refinement, we compute an importance score for every Gaussian primitive and remove those with minimal rendering contribution, thereby accelerating convergence and reducing redundant computation and memory overhead. To further enhance surface reconstruction quality, we also incorporate normal priors from a pretrained model. Finally, even under memory-constrained conditions, our method enables high-quality, high-resolution 3D scene reconstruction. Extensive experiments on three public benchmarks demonstrate that our approach achieves state-of-the-art performance in high-resolution novel view synthesis (NVS) and surface reconstruction tasks.
\end{abstract}



